%======================================================================
\section{Related work}
\label{sec:related}

\paragraph{First- and second-order removal estimates.}
\citet{mozer1989skeletonization} estimated the relevance of a unit from the derivative of the network's error with respect to a gain on the unit's output, the construction we evaluate here with silencing strength as the gain. Optimal brain damage and optimal brain surgeon used the diagonal or the full Hessian of the loss instead~\citep{lecun1990optimal,hassibi1993second}, and \citet{molchanov2017pruning,molchanov2019importance} ranked convolutional filters and neurons for pruning with first- and second-order Taylor expansions of the loss. In pruning the estimate has to rank components; our question is whether it also gives the size and sign of each effect, and if not, why not.

\paragraph{Attribution patching.}
In language models, attribution patching approximates activation patching to first order~\citep{nanda2023attribution}. Edge attribution patching recovered circuits better than earlier automated methods~\citep{syed2024attribution}. \citet{kramar2024atp} identified two kinds of failure that cause false negatives, saturated attention and cancellation, and addressed the first by computing the change of the attention softmax exactly instead of linearising it when queries and keys are patched. That repair is of the same kind as the estimator that works best in Study 2, which keeps the inputs to a normalising step linear and applies the normalisation exactly. \citet{hanna2024faith} found that circuits chosen with integrated gradients~\citep{sundararajan2017axiomatic} are more faithful than those chosen with plain attribution patching. Attribution patching can lose the size of the effects almost entirely: for the outputs of the MLP layers of GPT-2 Large in the indirect-object identification task, \citet{jafari2025relp} found a Pearson correlation of 0.006 between attribution and activation patching, and proposed replacing the local gradients with coefficients from layer-wise relevance propagation. \citet{zhang2026attribution} traced the dominant error of attribution patching to nonlinearities downstream of the patched component and proposed a second-order correction computed with Hessian-vector products, and a workflow that screens with first-order estimates and spends further computation only on the components flagged as unreliable. In the FlyVis models that correction makes the estimates worse, because the curvature it adds comes mostly from a ratio whose Taylor series diverges for large perturbations, and because it cannot see the network's threshold crossings. Our setting differs from these studies in three ways: the network is recurrent and the output is a peak-based tuning index, the intervention silences a whole cell type, and the exact effect of every intervention is known, so effect-size errors can be measured directly.

\paragraph{Baselines and attribution in the FlyVis models.}
The value that replaces a component's output (zero, a mean, or a resampled value) changes what patching experiments find~\citep{heimersheim2024patching,miller2024metrics}, and \citet{heimersheim2024patching} discuss how the choice of metric, for example a logit difference against a probability, affects the conclusions. In work concurrent with ours, \citet{dhawan2026flyhack} scored integrated gradients and edge attribution patching against exhaustive interventions on each type-to-type connection in ten of the FlyVis models, with decoded optic flow as the output. Both methods ranked the important connections highly whatever value replaced a removed connection, but they ranked the sizes of the effects well only when built for the same intervention as the reference: against the effects of mean replacement, edge attribution patching reached a Spearman correlation of 0.872 when built for mean replacement and 0.126 when built for zero replacement. That study scores ranks and classification. In Study 1 the estimate and the reference share one intervention, so the effect-size errors come from nonlinearity rather than from a mismatch of baselines, and we measure absolute and signed errors against predicting zero change, how they change with silencing strength, and the cost of each method. In Study 2, changing the shared baseline to each cell's mean output reduces the error but does not remove it, and a mismatched derivative loses both the effect sizes and the screening advantage.

\paragraph{Connectome-constrained models of the fly visual system.}
\citet{lappalainen2024connectome} trained the FlyVis ensemble and described groups of models with different T4c tuning. During peer review they simulated silencing of Mi1, Tm3, Mi4 and Mi9 by clamping each type's outputs to zero, averaged the T4c responses over the models with correct T4c tuning, and reported qualitative agreement with the experiments of \citet{strother2017emergence} for most effects~\citep{lappalainen2024review}. \citet{ryan2026silencing} silenced cell types in the same \NModels{} models and compared the effects with published results of genetic silencing, under a protocol fixed before any outcome was coded: the ensemble's predictions agreed with the experiments better than chance but no better than a baseline built from synapse counts alone, which, like our anatomical path weights, follows paths of up to three synaptic steps back from the readout. Later studies in the same series simulated every single cell type and every pair, and samples of up to five types together, in every model~\citep{ryan2026tests,ryan2026budget}; there, a linear model of the signed connectome ranked the magnitudes of the interactions between silenced types but predicted their signs at or below chance~\citep{ryan2026budget}, a split between locating effects and estimating them like the one we find for first-order estimates. \citet{dhiman2026scale} silenced T4 and T5 in a FlyVis model and found that a refitted linear decoder still recovered the direction of motion. \citet{karaneen2026reproducibility} retrained ensembles of such models and found that selecting models by task error does not reliably identify the solutions that match measured tuning. \citet{pospisil2024effectome} proposed estimating a linear causal model of the fly brain from optogenetic perturbations, with the connectome as a prior; our anatomical path weights use the connectome alone in a related way. \citet{xie2025visual} predicted the visual responses of neurons across the fly brain by aggregating paths through powers of the connectome's adjacency matrix, a construction close to our path weights.

\paragraph{Biological silencing.}
Genetic silencing experiments characterised the inputs of the ON motion detector T4, Mi1, Tm3, Mi4 and Mi9~\citep{ammer2015functional,strother2017emergence}, and of the OFF motion detector T5, Tm1, Tm2, Tm4 and Tm9~\citep{serbe2016comprehensive}; \citet{borst2020fly} review the circuit. Inhibition relayed by CT1 shapes the tuning of T5~\citep{braun2023disynaptic}, and the feedback neurons C2 and C3 take part in motion processing~\citep{henning2026inhibitory}. These experiments measure other quantities in other cells, by calcium imaging or electrophysiology, after silencing that lasts much longer than ours, so we use them as context only (\cref{app:biology}). Elsewhere in the fly brain, silencing in a connectome model has been tested against experiments: in a leaky integrate-and-fire model of the whole central brain, \citet{shiu2024drosophila} silenced neurons one at a time and tested several of the predicted effects on feeding by optogenetic silencing in flies.
